% HamSCI oral presentation template: a worked example that follows the slide order of the
% PowerPoint template and shows each feature of the Beamer theme. Copy this file to start a talk.
% Build from this directory with latexmk; latexmkrc points TEXINPUTS at ../beamer.
\documentclass[aspectratio=169,10pt]{beamer}
\usetheme{HamSCI}
\usepackage{booktabs}
\usepackage{pgfplots}
\pgfplotsset{compat=1.18}

\title{Presentation Title}
\subtitle{Optional subtitle}
% \texorpdfstring keeps the affiliation numbers out of the PDF's author metadata.
\author{\textbf{Author One KC4AAA}\texorpdfstring{\textsuperscript{1,2}}{}, Author Two KC4USV\texorpdfstring{\textsuperscript{1}}{}}
\institute{\textsuperscript{1}HamSCI\\\textsuperscript{2}Institution Two}
\date{Meeting name, City, Day Month Year}
\hamscicontact{email@domain.edu}

\begin{document}

\hamscititleframe


% ---------------------------------------------------------------------------------------------
\begin{frame}{Introduction}
\hamscitakeaway{Open every slide with one line that says what the audience should remember}
\begin{itemize}
  \item Then support it with short lines, a table, or a figure, not paragraphs.
  \item Keep fonts large enough to read from the back of the room.
  \item Use American English throughout: color, center, modeled, labeled.
\end{itemize}
\end{frame}


% ---------------------------------------------------------------------------------------------
\begin{frame}{Methodology}
\hamscitakeaway{Tables carry numbers better than sentences do}
\small
\begin{tabular}{@{}lll@{}}
\toprule
Instrument & Network & What it measures \\
\midrule
HF receivers & PSKReporter, RBN, WSPRNet & who heard whom, signal-to-noise ratio \\
PSWS Grape & HamSCI PSWS & Doppler shift of time-standard stations \\
Ionosonde & GIRO & electron density profile \\
\bottomrule
\end{tabular}\par\vskip10pt
\footnotesize Details and caveats go in a footnote-sized line below the table.
\end{frame}


% ---------------------------------------------------------------------------------------------
\begin{frame}{Results}
\begin{columns}[c]
\begin{column}{0.64\textwidth}
\begin{tikzpicture}
\begin{axis}[width=\linewidth, height=5.6cm, xlabel={UT (hours)}, ylabel={Value (units)},
             xmin=0, xmax=24, grid=major, grid style={gray!25}, tick label style={font=\small},
             label style={font=\small}]
  \addplot[very thick, color=hamsciaccent, domain=0:24, samples=97]
    {2 + sin(deg(2*pi*x/24)) + 0.3*sin(deg(2*pi*x/3))};
\end{axis}
\end{tikzpicture}
\end{column}
\begin{column}{0.34\textwidth}
\hamscitakeaway{Put the figure first, and the text in the whitespace beside it}
\small Size the figure to fill the slide, then fit the words around it.
\end{column}
\end{columns}
\hamscisource{Figure: \href{https://github.com/HamSCI/presentation_templates}{link to the PNG}
  \textbullet\ code: \href{https://github.com/HamSCI/presentation_templates}{link to the script, pinned to a commit}}
\end{frame}


% ---------------------------------------------------------------------------------------------
\begin{frame}{Discussion}
\hamscitakeaway{Two columns let you compare without crowding}
\begin{columns}[t]
\begin{column}{0.48\textwidth}
\textbf{What the data show}
\begin{itemize}
  \item First point
  \item Second point
\end{itemize}
\end{column}
\begin{column}{0.48\textwidth}
\textbf{What remains open}
\begin{itemize}
  \item First question
  \item Second question
\end{itemize}
\end{column}
\end{columns}
\end{frame}


% ---------------------------------------------------------------------------------------------
\begin{frame}{Summary}
\hamscitakeaway{Three to five numbered points, each one line}
\begin{enumerate}
  \item \textbf{Observed}: the main result, with its number.
  \item \textbf{Interpretation}: what it means.
  \item \textbf{Next}: what comes next.
\end{enumerate}
\end{frame}


% ---------------------------------------------------------------------------------------------
\hamsciacknowledgments[This work is supported by \textit{list the awards that funded this work}.]


% ---------------------------------------------------------------------------------------------
% Include this slide whenever AI tools helped make the talk or the work behind it, and keep it
% accurate. The project's AI usage log is the record it summarizes.
\begin{frame}{Disclosure of AI use}
\small
\begin{itemize}
  \item \textbf{Tool}: name and exact model version, for example Claude (Anthropic),
    \texttt{claude-opus-5-5}, via Claude Code.
  \item \textbf{When}: the dates of the AI-assisted sessions.
  \item \textbf{How it was used}: concretely, for example analysis code, figures, or drafting of
    these slides.
  \item \textbf{Human oversight}: the authors decided the research questions, methods and
    interpretation, reviewed all AI-assisted output, and are responsible for all content.
  \item \textbf{Record}: where the session log is kept, for example \texttt{ai/ai\_usage\_log.md}
    in the project repository.
\end{itemize}
\end{frame}


% ---------------------------------------------------------------------------------------------
\hamscisectionframe{Thank you!}

\end{document}
